% [R1-BASE] Restored from 5a8b0be (42-page source), then display-only outline migration.
% Latest earlier prose preserved in revisions/saved_ad78fcd. R2 now revises only the active abstract.
% See REVISION_STEP1.md for routing reasons and deferred edits. No deletion approved.
%% PeerJ Computer Science submission manuscript.
%% The lineno option is required for peer-review submissions; remove it only
%% for a camera-ready version when requested by PeerJ.
\documentclass[fleqn,12pt,lineno]{wlpeerj}

%% Packages needed by the existing manuscript content. The PeerJ class already
%% loads graphicx, booktabs, xcolor, caption, natbib, and microtype.
\usepackage{color,xspace}
\usepackage{array}
\usepackage{ragged2e}
% [R12-L01] Ordinary tabularx avoids ltablex incrementing table counters inside floats.
\usepackage{tabularx}
% [R18] Breakable handbook tables; ordinary tabularx remains unchanged elsewhere.
\usepackage{longtable}
\usepackage{makecell}
\usepackage{placeins}
\usepackage{needspace}
\usepackage{adjustbox}
\usepackage[hidelinks]{hyperref}
\usepackage{float}
\usepackage{pdfpages}
\usepackage{amssymb}
\usepackage{colortbl}
\usepackage{multirow}
\usepackage{rotating}
\usepackage{fontspec}
\usepackage{xeCJK}
\usepackage{enumitem}
% PeerJ CS initial-submission formatting, checked 2026-09-18.
\geometry{letterpaper,left=2.5cm,right=2.5cm,top=2.5cm,bottom=2.5cm,headheight=14pt}
\setmainfont{TeX Gyre Termes}
\usepackage{etoolbox}
\makeatletter
\patchcmd{\@maketitle}{\sffamily\small\vskip1ex}{\rmfamily\normalsize\RaggedRight\vskip1ex}{}{\PackageError{cnss}{Abstract font patch failed}{Check wlpeerj title definition.}}
\makeatother
\hypersetup{pdfsubject={PeerJ Computer Science research article}}
\lfoot{\footnotesize\textit{PeerJ Computer Science}}

% [R44] Bundled CJK fonts produce embedded Unicode maps with XeLaTeX.
% Keep Chinese as selectable text; do not depend on a reader's Adobe-GB1 maps.
\setCJKmainfont[
  Path=fonts/noto-cjk/,
  BoldFont=NotoSerifCJKsc-Bold.otf,
  ItalicFont=NotoSerifCJKsc-Regular.otf,
  ItalicFeatures={FakeSlant=0.2}
]{NotoSerifCJKsc-Regular.otf}
\setCJKsansfont[Path=fonts/noto-cjk/]{NotoSansCJKsc-Regular.otf}
\setCJKmonofont[Path=fonts/noto-cjk/]{NotoSansCJKsc-Regular.otf}

\raggedbottom
\setlength{\emergencystretch}{1em}
\widowpenalty=10000
\clubpenalty=10000
% [R12-L02] Let readable figures share text pages rather than forming sparse float pages.
\renewcommand{\topfraction}{0.9}
\renewcommand{\bottomfraction}{0.8}
\renewcommand{\textfraction}{0.08}
\renewcommand{\floatpagefraction}{0.8}
%% The official PeerJ template uses unnumbered section headings. Keep the
%% numbered source files unchanged and apply that presentation globally here.
\setcounter{secnumdepth}{-1}






%% Column types used in the manuscript tables.
\newcolumntype{L}{>{\RaggedRight\arraybackslash}X}
\newcolumntype{P}[1]{>{\RaggedRight\arraybackslash}p{#1}}

%% Language helpers retained from the source manuscript.
\newcommand{\zh}[1]{#1}
\newcommand{\en}[1]{#1}

%% AML laboratory review mode. Wenlin's comments and all revision wrappers
%% remain in the editor; keep this false for a clean PeerJ-facing PDF and set
%% it true only when collaborators need a visible annotated review copy.
\newif\ifamlreview
\amlreviewfalse

%% Supervisor comments remain verbatim at their original source locations.
\newif\ifshowwlcomments
\ifamlreview
  \showwlcommentstrue
\else
  \showwlcommentsfalse
\fi
\ifshowwlcomments
  \newcommand{\wl}[1]{\textcolor{purple}{\textbf{wl:} #1}}
\else
  \newcommand{\wl}[1]{}
\fi

%% AML-lab revision convention: substantial revised passages can be shown in
%% blue during internal review.  The master switch above controls both the
%% revision highlighting and collaborator-comment display.
\newif\ifshowrevisions
\ifamlreview
  \showrevisionstrue
\else
  \showrevisionsfalse
\fi
\newcommand{\rev}[1]{\ifshowrevisions{\color{blue}#1}\else#1\fi}
\newcommand{\zxy}[1]{{\color{blue}[#1 -- Xy]}}
\newcommand{\zc}{{\color{blue}[cite]}\xspace}
\newcommand{\etal}{\emph{et al.}\xspace}
\newcommand{\eg}{\emph{e.g.,}\xspace}
\newcommand{\ie}{\emph{i.e.,}\xspace}
\newcommand{\etc}{\emph{etc.}\xspace}

%% Reusable figure floats are defined here and invoked only at their
%% section-specific locations below.
\input{tex/skillspan_style_figures}
\input{tex/stage2_result_values}

\title{Chinese-SkillSpan: A Versioned Resource for Competency Span Extraction from Chinese Job Advertisements}

\author[1,2]{Guojing Li\textsuperscript{\textdagger}}
\author[2]{Zichuan Fu\textsuperscript{\textdagger}}
\author[2]{Junyi Li}
\author[2]{Wenlin Zhang}
\author[2]{Kaifeng Guo}
\author[2]{Jinning Yang}
\author[3]{Xinyang Wu}
\author[2]{Jingtong Gao}
\author[2]{Xiangyu Zhao}

\affil[1]{Renmin University of China, Beijing, China}
\affil[2]{City University of Hong Kong, Kowloon, Hong Kong SAR, China}
\affil[3]{Wuhan University, Wuhan, Hubei, China}

\corrauthor{Xiangyu Zhao}{xianzhao@cityu.edu.hk}

%% PeerJ submission keywords (enter these in the submission system):
%% Chinese JobSkillNER; competency span extraction; skill extraction;
%% Chinese job advertisements; benchmark dataset; span-level annotation;
%% domain-adaptive pretraining; large language models; ESCO-informed taxonomy.

\begin{abstract}
\par\noindent\rev{\textbf{Background.} Chinese job advertisements often express competencies through coordinated phrases and context-dependent terms. Extracting these mentions requires consistent decisions about their boundaries and types.}
\par\noindent\rev{\textbf{Methods.} We developed Chinese-SkillSpan from 22,840 sentences across four recruitment collections. An ESCO-informed handbook defines character spans for language, knowledge, occupational skills, and transversal competences. LLM-generated labels support training and evaluation, while a fixed 150-sentence human reference supports diagnostic assessment. The reference also contributed to guideline development. A separate three-coder study assessed final-handbook agreement on 50 sentences; historical coding and assisted review provide additional quality checks. We evaluated six inference configurations, supervised adaptation, and multilingual encoder baselines.}
\par\noindent\rev{\textbf{Results.} Under fixed inference rules, Qwen adaptation increased typed exact F1 from 0.361 to $0.540\pm0.035$ across three training seeds. Chinese-supervised XLM-R and ESCOXLM-R achieved 0.522 and 0.538; their paired interval included zero. The new three-coder study yielded mean typed exact span F1 of 0.663 (95\% CI [0.562, 0.759]) and character-level $\alpha=0.817$, revealing remaining span disagreements.}
\par\noindent\rev{\textbf{Conclusions.} Chinese-SkillSpan makes boundary and type decisions explicit and enables comparison of generated supervision with human annotations. The present evaluation concerns a reference used during development. Independent coding under the final handbook identifies remaining annotation ambiguity; generalization beyond the development-used reference remains unestablished.}
\end{abstract}

% External benchmark results are integrated into the main experiments and Appendix E.

\begin{document}
\RaggedRight

\raggedbottom
\maketitle
\thispagestyle{empty}

\begingroup
\renewcommand\thefootnote{\textdagger}
\footnotetext{Guojing Li and Zichuan Fu contributed equally to this work.}
\endgroup

% [R4-INTRO] Original introduction retained; current narrative follows Outline V2.
\input{1Introduction}
\input{5Relatedwork}
\input{2Framework}
\input{3Experiments}
\section{Limitations and Responsible Use}
\label{sec:limitations}

\textbf{Reference design and coverage.} The human reference was selected for calibration and challenge analysis and was repeatedly examined during guideline development. Exclusion from gradients and checkpoint selection does not remove that dependence. Its small, uneven type distribution limits generalization, and keyword-prioritized expansion does not represent the recruitment population. The new final-handbook study supplies independent three-coder evidence, but mean typed span F1 is 0.663 and L has only two spans. Its 50-sentence teacher-pool sample does not establish corpus-wide reliability. Historical coding and assisted review use different conditions. The annotations describe advertised requirements, not individual applicants' abilities.

\textbf{Data access and reproducibility.} Public materials allow inspection of labels, scoring of saved predictions, and comparison of the released historical coding and review layers. They do not include the complete final expanded-pool training inputs; the annotation release preserves analytical fields but omits administrative account metadata. Some collection records, historical implementation details, and label-version compatibility checks are incomplete. These gaps limit full retraining and verification. Data Availability specifies component-level access and reuse terms; applications also require source permissions, privacy safeguards, human oversight, and external validation.

\textbf{Experimental comparisons.} Sentence-level splits, residual JobBERT training/development text matches, and the inherited CRF's incomplete selection history limit claims about independent generalization. The supervision studies differ in what they change, and expanded-pool runs use single seeds. Inference budgets vary, while some proxy-reported model identities remain unverified. External datasets also differ in labels, splits, scorers, and selection procedures. Their bootstrap intervals condition on the trained seeds and do not address advertisement clustering or multiple comparisons. Chinese encoder predictions have been rescored, but their small paired difference does not establish an initialization advantage. These comparisons describe the tested settings rather than a general ranking of models or languages.

\section{Conclusion}
\label{sec:conclusion}

Chinese-SkillSpan specifies how to delimit and type competency mentions in Chinese recruitment text, with explicit treatment of coordination, contextual types, and negative cases. Its contribution is a versioned resource and an inspectable annotation procedure. The new three-coder study measures final-handbook consistency separately from historical coding and assisted review. Its remaining boundary and type disagreements show where adjudication and rule clarification are still needed.

On the development-used human reference, Qwen adaptation improved end-to-end extraction under fixed inference rules. The additional Chinese-supervised encoders provide reproducible alternatives, but the paired interval does not establish an ESCO initialization advantage. Boundary-sensitive scores and the type, length, and empty-sentence diagnostics identify where further assessment is needed. These findings support diagnostic use of the released resource. The new agreement results provide evidence about the annotation process, but do not establish model generalization to new advertisement documents or resolve all historical label-version differences.


% [R4-LAYOUT] Keep each repository URL intact without forcing an extra page.
\section{Data Availability}
\begin{sloppypar}
\RaggedRight % [R4-LAYOUT] Avoid stretched word spacing around intact URLs.
% [R39-DAS] Two-paragraph access statement; detailed version/model inventory remains in Appendix D.
The Chinese-SkillSpan dataset, annotation guidelines, predefined data splits, and documentation are available in the \href{https://doi.org/10.5281/zenodo.22698504}{Zenodo archive (v0.1.3; DOI: 10.5281/zenodo.22698504)}, under \href{https://doi.org/10.5281/zenodo.22288337}{concept DOI 10.5281/zenodo.22288337}. This version includes the 150-sentence human reference and the Silver-plus B2 training/development files. Later expanded-Silver experiments and QA150 review materials are outside this dataset archive. Pseudonymized historical blind-coding and dual-reviewer exports are separately available in the \href{https://github.com/AlfredJamesLi/chinese-skillspan-benchmark/blob/main/reproduction/evidence_review_20260925/annotation_layers/README.md}{annotation-layer release}; original text and annotation offsets are preserved. The new three-coder study is released separately with \href{https://github.com/AlfredJamesLi/chinese-skillspan-benchmark/tree/main/reproduction/agreement/finalguide_abc_20261002}{pseudonymized frozen annotations, source identifiers, agreement code, and results}; it is not included in the original Zenodo dataset archive.

Qwen LoRA adapters, frozen evaluation predictions, configurations, and scoring tools are available in a separate \href{https://doi.org/10.5281/zenodo.22851581}{model-reproduction archive (DOI: 10.5281/zenodo.22851581)}. It includes the three shared-guideline B2 adapters and selected expanded-Silver checkpoints. The complete texts and labels for the final 9,540- and 9,646-record pools, including their training, validation, and test partitions, remain outside the public archives; the expanded-study package supports inference and rescoring, but not complete retraining from public inputs alone. The \href{https://github.com/AlfredJamesLi/chinese-skillspan-benchmark}{project repository} provides source code, preprocessing scripts, evaluation tools, and a \href{https://github.com/AlfredJamesLi/chinese-skillspan-benchmark/blob/main/REPRODUCIBILITY.md}{reproduction guide}.

The six Chinese-supervised XLM-R/ESCOXLM-R checkpoints are publicly available through the \href{https://github.com/AlfredJamesLi/chinese-skillspan-benchmark/tree/0c73207c0f9b00fa87f385b78bcb7342bd00bfcb/results_snapshots/table11_evidence_20260924/HF_RELEASE.md}{model index}. Their predictions, executed code, configurations, and verification records are preserved in the \href{https://doi.org/10.5281/zenodo.22937245}{Chinese encoder audit archive (DOI: 10.5281/zenodo.22937245)}; model weights are hosted on Hugging Face. The \href{https://github.com/AlfredJamesLi/chinese-skillspan-benchmark/blob/main/reproduction/evidence_review_20260925/README.md}{evidence index} distinguishes locally executed rescoring from server-reported checkpoint inference and records the outstanding annotation-provenance checks.

The \href{https://huggingface.co/AlfredJames/jobbert-zh}{JobBERT-zh initialization} provides the domain-adapted encoder and inherited CRF. The \href{https://huggingface.co/AlfredJames/jobbert-zh-v6a}{JobBERT-zh B2 continuation} provides the continued checkpoints; its default checkpoint is seed 42, rather than the three-seed mean. Reuse terms vary by component. Software and documentation explicitly covered by Apache-2.0 retain that licence; author-controlled annotations and research materials permit academic research and peer review with attribution. Third-party recruitment text is excluded from that author-granted permission. Model checkpoints and adapters follow their accompanying notices and applicable base-model licences. The \href{https://github.com/AlfredJamesLi/chinese-skillspan-benchmark/blob/main/DATA_AVAILABILITY.md}{data-access guide} specifies the component-level terms and available materials. Earlier archive versions are identified in \hyperref[app:d]{Appendix D}. Table~\ref{tab:component-access} distinguishes the available objects and reproduction levels. External benchmark protocols and their versioned result ledger are documented in \hyperref[app:e]{Appendix E}.
\end{sloppypar}

% [R16] Four reader-facing appendices replace chronological/provenance-heavy routing.
% Original inputs remain in the project and the R15 git snapshot; no assets deleted.
\input{tex/appendix_R16}
\input{tex/external_benchmarks/paper_appendix}
\FloatBarrier

% Funding and references use the submission-wide geometry.
% Use the same 2.5 cm margins for declarations and references.
\section{Funding}
This research was supported by the National Social Science Fund of China under
the project ``Research on the Interaction Mechanism between Enterprise Employees'
Social Capital and Human Capital in the Intelligent Era and Its Impact on Career
Success'' (Grant No.~21BGL142).
The funder had no role in study design, data collection and analysis, manuscript preparation, or the decision to submit the article for publication.

\section{Competing Interests}
The authors declare that they have no competing interests.

\section{Acknowledgments}
OpenAI Codex assisted language polishing, LaTeX formatting, and code preparation; OpenAI's image-generation tool (\texttt{image\_gen}) assisted the schematic artwork in Figure~\ref{fig:macro-micro} and its revision under author direction. These tools were used in September 2026; their model-version identifiers were not retained. The authors take responsibility for the manuscript, references, code, and figures, including verification of AI-assisted content. Language models used in annotation and experiments are documented in the methods and reproduction materials.

% Funding role and competing interests confirmed by the author on 2026-09-18.

%% Keep ordinary widow/club penalties in the reference list.
\begingroup
\widowpenalty=3000
\clubpenalty=3000
\bibliography{9Reference}
\endgroup


\end{document}
